\appendix
\renewcommand{\thesection}{\Alph{section}}
\section{Detailed circuit and execution protocol}\label{app:protocol}
\subsection{Circuit families and observables}
All circuits start from ideal encoded Bell pairs with noiseless references. A perfect final extraction closes the detector record, and final commuting Bell products monitor both logical Pauli sectors. For $B$ blocks encoding $k_q$ logical qubits each, there are $N_{\rm obs}=2Bk_q$ monitored outputs. The Bell-reference construction follows Appendix A.7 of Ref.~\cite{tour}.

Surface-code circuits use rotated $[[d^2,1,d]]$ codes with $d\in\{3,5,7\}$, $B=4$, $d$ noisy extraction rounds, and $p\in\{0.001,0.003,0.006\}$. Each block uses $d^2-1$ extraction ancillas and one noiseless reference in addition to its data qubits. Each noisy round begins with a nearest-neighbor transversal CNOT matching. Matchings alternate, and directions reverse every two rounds. Ancilla resets have $X$-flip noise, ancilla $H$ gates and all CNOT gates have depolarizing noise, and readout bits flip with probability $p$. Idle noise is omitted.

BB72 and BB144 use the $[[72,12,6]]$ and $[[144,12,12]]$ constructions from Ref.~\cite{bbpre}. Four blocks undergo six noisy slices at $p=0.005$ or 0.015. Inter-block CNOTs and commuting stabilizer-product measurements are ideal except for independent readout flips; each slice applies single-qubit depolarization to data. These controls use a two-sector phenomenological noise model. Each monitors 96 Bell observables. Additional eight-block surface cases at $p=0.003$, $d=3,5$, test multiple seams and support the engineering audit. Table~\ref{tab:grid} lists the circuit and sampling parameters. The distance-seven cases have limited statistical power.

Stim~\cite{stim} generates an undecomposed detector error model with no approximation of disjoint errors. We sample the model with a fixed seed and verify $De=s$ and $Le=\ell$ for every shot. Noiseless tests and 54 injected logical $X/Z$ checks validate the Bell-reference observables. No approximate graphlike projection is used to make the correlated hypergraph acceptable to a matching decoder.

\begin{table*}[t]\centering\small
\caption{Circuit parameters and sample counts for the 15 configurations. Each of BB72 and BB144 uses the listed shot counts. Dashes indicate noise settings outside the grid.}\label{tab:grid}
\begin{tabular}{lrrrrrrrr}\toprule
&&&\multicolumn{5}{c}{Shots at circuit noise parameter $p$}\\
Code&Blocks&Noisy rounds&0.001&0.003&0.005&0.006&0.015\\\midrule
Surface $d=3$&4&3&512&512&--&512&--\\
Surface $d=5$&4&5&256&1,024&--&256&--\\
Surface $d=7$&4&7&64&64&--&64&--\\
BB72, BB144&4&6&--&--&128&--&128\\
Surface $d=3$&8&3&--&256&--&--&--\\
Surface $d=5$&8&5&--&64&--&--&--\\\bottomrule
\end{tabular}\end{table*}

\subsection{Decoder stacks and local row selection}
The four engine settings are BF, BP+LSD0, BP+LSD2, and BP+OSD0. BF uses inversion and one-bit growth; LSD2 uses combination-sweep order two, and OSD0 uses order-zero ordered-statistics postprocessing~\cite{lsd,osd,bfdocs,lsddocs,osddocs}. All share 30 min-sum BP iterations, scaling 0.625, parallel message scheduling, and one native CPU thread. We use Stim 1.16.0 and LDPC 2.4.1. The engine grid covers three postprocessing families and two LSD orders. It does not establish optimality among all correlated decoders.

Each engine is crossed with full joint decoding, direct cutting, and radius-one repair. BF additionally tests halos of one and two blocks, radius-two repair, and rank-screened radius-one repair, giving 16 stacks. The core width is two. In four-block models, radius two unfreezes all nonzero detector columns and therefore performs a full-width re-decode filtered against the cut incumbent. In eight-block models, its four-owner-block bands are proper subsets. For a halo window, the local row set includes every touched row with no remaining outside-window variable, together with all core rows. A failed local solve or unresolved residual triggers a full joint call using the same engine; fallback frequency and its full latency are recorded. Seam replacement uses all touched rows, fixed exterior variables, and a $10^{-10}$ acceptance tolerance. Every returned correction is audited against the complete $D$ matrix.

The frozen resource allocation includes LSD order two and OSD order zero. Higher-order OSD is outside the main grid, so the results do not rank boundary repair against that baseline. The protocol records an initial timing pilot as the allocation rationale; its durations are not used as quantitative evidence in the revised comparisons. Matching-based circuit representations~\cite{fast} and other modern BP variants remain relevant external comparisons.

All conditional targets in the study, including warmups, pass the exact column-space membership check. The guard therefore rejects zero native calls in this dataset. Long durations are retained as measured decoding cost and are not classified as evidence of infeasibility.

\subsection{Sampling, statistics, and multiplicity}
We set the configurations, seeds, shot counts, engine settings, and nine planned contrasts per model before initial sampling. During execution, problematic native calls motivated an implementation amendment: an exact binary column-space guard rejects an infeasible local target before calling the native engine. It implements the existing failed-local-solve/fallback rule and leaves feasible inputs unchanged. All configurations were rerun with this guard, whose cost is included in timing. We made the amendment after inspecting partial outcomes. Sample counts, seeds, hypotheses, and engine settings were unchanged. Methods share the same samples and run in seeded random order after three excluded warmups. There is no outcome-dependent stopping. The 15 models contribute 135 planned paired tests to one Holm family~\cite{holm}.

We use 95\% Wilson intervals for marginal failure rates~\cite{wilson}. For a paired contrast, $n_{10}$ counts baseline failures fixed by the comparison method, and $n_{01}$ counts new failures introduced by that method. The effect is $(n_{10}-n_{01})/N$. We use the exact conditional binomial version of the paired McNemar test~\cite{mcnemar} and conservative 95\% paired-effect intervals obtained by subtracting two 97.5\% Clopper--Pearson intervals~\cite{cp}; the union bound supplies simultaneous coverage for the two discordance probabilities. These effect intervals are not adjusted across all contrasts. Appendix~\ref{app:statistics} gives the formulas.

\subsection{Timing and memory estimands}
Main timings include decoding, wrapper operations, consistency checks internal to the policy, and repair-event construction; compilation, sampling, file writes, and the final outcome audit are excluded. Model-level simulation jobs run with up to three concurrent processes on one shared host, each native decoder using one thread; the principal distance-five case was completed with one model job active. These timings describe that execution environment. The separate engineering audit serializes the compared method runs and provides the controlled serial/parallel resource comparison.

The engineering audit runs three fresh-process replicates of each method on 64 common eight-block samples, with seeded randomized method order. It compares event-free serial and two-worker process-parallel repair with identical checks and acceptance arithmetic. Initialization is reported separately; the processes and local solvers remain warm during timed calls. Memory is instrumented in a separate pass over the first eight audited inputs. Per-process OS high-water RSS, covering the preceding timed calls as well, and 2 ms samples of aggregate process-tree RSS and proportional set size (PSS) are retained. Sampled aggregate peaks can miss short transients and uninstrumented inputs; PSS avoids counting shared pages in full for every process.

\section{Additional surface-code and rank results}\label{app:surface}
Figure~\ref{fig:surface} shows the original surface-code sweep at distances three, five, and seven. Within the 15 configurations, radius-one BF repair yields lower observed failure counts than BF cutting in 13 cases and ties in two. All paired contrasts remain in the 135-test correction family. Figure~\ref{fig:rankcurve} traces logical reach as fault columns are added to a repair band. The original figures are retained; the controlled backend study additionally reuses the stored sample archives and records their hashes.

\begin{figure*}[t]\centering\includegraphics[width=\textwidth]{figures/fig5.pdf}
\caption{Surface-code sweep at distances 3, 5, and 7. Methods use the same fault samples within each configuration. Error bars show marginal 95\% Wilson intervals.}\label{fig:surface}\end{figure*}
\begin{figure*}[t]\centering\includegraphics[width=\textwidth]{figures/fig6.pdf}
\caption{Logical reach as the selected fault-column set grows. Columns enter in fixed prior order. The rank calculation uses the circuit matrices without shot-level logical truth.}\label{fig:rankcurve}\end{figure*}

\section{Complete engineering measurements}\label{app:engineering}
The audit serializes the measured method runs. Each method uses three fresh processes on 64 shared samples. Mean and p95 pool those durations; memory entries are maxima of separately instrumented eight-input passes. The main-text speed ratio excludes parallel-worker startup. Table~\ref{tab:engineering} includes both measured distances.
\begin{table*}[t]\centering\small
\caption{Warm-call latency and process-tree memory for the eight-block surface-code circuits. Timing and memory were measured in separate passes.}\label{tab:engineering}
\begin{tabular}{lrrrrrr}\toprule
&\multicolumn{3}{c}{$d=3$}&\multicolumn{3}{c}{$d=5$}\\
Stack&Mean ms&p95 ms&PSS MiB&Mean ms&p95 ms&PSS MiB\\\midrule
BF joint&2.58&5.99&66.6&68.78&149.46&74.6\\
BF cut&2.57&4.86&68.9&35.44&77.20&82.3\\
BF repair1&3.53&6.57&70.3&60.00&137.54&88.1\\
BF parallel1&4.57&8.19&193.0&49.71&122.24&233.1\\
LSD2 joint&2.57&6.73&66.7&54.39&110.84&74.7\\
LSD2 repair1&3.39&6.19&70.0&61.77&155.71&87.7\\
OSD0 joint&2.50&5.26&67.4&51.34&74.21&79.8\\
OSD0 repair1&3.22&6.68&72.3&51.86&73.74&101.6\\\bottomrule
\end{tabular}\end{table*}

\section{Ownership-based initial partition pass}\label{app:initial}
For a core width $q$, process core block sets
\[C_t=\{t,\ldots,\min(t+q-1,B)\},\quad t=1,1+q,\ldots,\ t\le B.\]
The halo width $h$ extends the current owner range to
\begin{equation}W_t=\{j:t\le\own(j)\le\min(t+q+h-1,B)\}.\label{eq:halo}\end{equation}
The inner decoder receives the residual syndrome restricted to a local row set $U_t$. The essential interface requirement is that $U_t$ include every detector row in the core. Additional rows provide context; their exact selection is an implementation choice. After a local candidate is obtained, only variables $K_t=\{j:\own(j)\in C_t\}$ are committed. Halo-owned estimates are discarded. A local candidate must close the core under the residual syndrome; a failure to do so must be flagged rather than silently committed.

\AlgBox{Algorithm A: ownership-based pre-commit overlap}{
\textbf{Input:} $D,s,\own,q,h$, row policy $U_t$, solver $\mathcal A$. Initialize $\hat e=0$, $s^{\rm res}=s$.
\begin{enumerate}
\item For $t=1,1+q,\ldots$, form $C_t,K_t,W_t,U_t$.
\item Compute $z\leftarrow\mathcal A(D_{U_t,W_t},s^{\rm res}_{U_t},p_{W_t})$. If no candidate closes every core row, return FAIL.
\item Set $\hat e_{K_t}\leftarrow z_{K_t}$ and $s^{\rm res}\leftarrow s^{\rm res}+D_{:,K_t}z_{K_t}$. Discard $z_{W_t\setminus K_t}$.
\item Return $\hat e$ only if $D\hat e=s$; otherwise return FAIL.
\end{enumerate}}
\begin{theorem}[Core closure]\label{thm:closure}
Suppose each commitment in Algorithm A closes the detector rows in its core. Then every closed row remains closed after all subsequent commitments. If all cores complete, the returned estimate satisfies $D\hat e=s$.
\end{theorem}
\begin{proof}
Consider a row $i$ in the current core, whose largest block label is $a$. Every column incident on $i$ has owner at most $\beta(i)\le a$. Earlier-owned columns have already been committed, and all remaining incident columns belong to the current core; later halo-owned columns contribute zero to row $i$. Thus a locally consistent candidate closes $i$ when its core variables are committed. Every subsequent committed column has owner greater than $a$ and is zero on row $i$ by Eq.~\eqref{eq:triangular}. Closure is therefore permanent. Induction over the cores closes every row.
\end{proof}
Core closure guarantees syndrome consistency when every core solve succeeds. A committed prefix can satisfy its current rows while making a later core infeasible. Appendix~\ref{app:examples} gives an example and motivates the joint fallback.

\paragraph{Numerical example of local submatrices.}
Return to Eq.~\eqref{eq:toy} and assign detector rows to blocks by $\beta(1)=1$, $\beta(2)=2$. The four column owners are $(1,1,2,1)$. For a one-block first core and one halo block, $C_1=\{1\}$, $K_1=\{1,2,4\}$, and $W_1=\{1,2,3,4\}$. With $U_1=\{1,2\}$, the local matrix $D_{U_1,W_1}$ is the full two-row matrix in Eq.~\eqref{eq:toy}; its solution uses all four columns as context, but only columns 1, 2, 4 are committed. These subscripts always select detector rows first and fault columns second.

The same columns illustrate conditioning and the rank identity. Select $J=\{1,2,4\}$, whose exterior is $\bar J=\{3\}$ and whose touched-row set is $R(J)=\{1,2\}$. Then
\[
D_J=\begin{pmatrix}1&1&1\\0&1&1\end{pmatrix},\quad D_{R(J),\bar J}=\begin{pmatrix}0\\1\end{pmatrix},
\]
\[L_{:,J}=\begin{pmatrix}1&0&1\end{pmatrix},\qquad t_J=\begin{pmatrix}s_1\\s_2+x_3\end{pmatrix}.\]
The exterior contributes a known offset to the second detector equation. For $s=(1,1)^T$ and $x_3=0$, both local candidates $(0,1,0)^T$ and $(0,0,1)^T$ satisfy $D_Jz=t_J$, but their logical actions are zero and one. Their difference $(0,1,1)^T$ lies in $\ker D_J$. Stacking the logical row below the detector rows gives
\[\begin{pmatrix}D_J\\L_{:,J}\end{pmatrix}=\begin{pmatrix}1&1&1\\0&1&1\\1&0&1\end{pmatrix},\quad\rho(J)=3-2=1.\]
Thus the stacked matrix counts constraints on both syndrome and logical action, whereas $D_J$ counts syndrome constraints alone. Their rank difference identifies the one logical bit that the local replacement can change. The acceptance test then chooses whether a proposed replacement lowers the prior-weight cost.

\section{BB matrices and logical basis validation}\label{app:bb}
\subsection{Translations and polynomial matrices}
Let $S_\nu$ be the cyclic shift acting on a unit row vector by $u_aS_\nu=u_{(a+1)\bmod\nu}$. For example,
\[S_3=\begin{pmatrix}0&1&0\\0&0&1\\1&0&0\end{pmatrix}.\]
On a periodic $\ell_x\times\ell_y$ lattice, define $T_x=S_{\ell_x}\otimes I_{\ell_y}$ and $T_y=I_{\ell_x}\otimes S_{\ell_y}$. A row indexed by site $(a,b)$ moves to $(a+1,b)$ under $T_x$ and to $(a,b+1)$ under $T_y$, with coordinates taken modulo the periods. Each binary monomial therefore specifies a translated support. With the code specification of Ref.~\cite{bbpre}, set
\[P=T_x^3+T_y+T_y^2,\qquad Q=T_y^3+T_x+T_x^2,\]
\[H_X=[P\mid Q],\qquad H_Z=[Q^T\mid P^T].\]
Here $P,Q$ each have size $\ell_x\ell_y\times\ell_x\ell_y$; concatenating them creates $n_q=2\ell_x\ell_y$ physical-qubit columns. Each row of $H_X$ gives the support of an $X$ stabilizer, and each row of $H_Z$ that of a $Z$ stabilizer. For BB72, $(\ell_x,\ell_y)=(6,6)$. Number sites within each half by $1+6a+b$. The first row of $P$ has ones at sites 19, 2, 3, and the first row of $Q$ at sites 4, 7, 13. Consequently the first row of $H_X$ specifies $X_2X_3X_{19}X_{40}X_{43}X_{49}$. This is one concrete stabilizer behind the compact polynomial notation. Since translations commute, $PQ=QP$, and
\[H_XH_Z^T=PQ+QP=0.\]
This binary product counts each pair's shared support modulo two, verifying that the $X$ and $Z$ stabilizers commute. The periods $(12,6)$ give BB144.

\subsection{Logical-operator support matrices}
The rows of $G_X$ and $G_Z$ are physical-qubit supports of logical Pauli operators. An $X$ support must be in $\ker H_Z$ to commute with every $Z$ check; adding a row combination of $H_X$ multiplies it by an $X$ stabilizer and preserves its logical action. Thus representatives are selected from $\ker H_Z/\operatorname{row}H_X$, and analogously $Z$ representatives from $\ker H_X/\operatorname{row}H_Z$. Kernels here are identified with their transposed row-vector supports. Canonical pairing requires
\[G_XH_Z^T=0,\qquad G_ZH_X^T=0,\qquad G_XG_Z^T=I_{k_q}.\]
For a small illustration, the $[[4,2,2]]$ CSS code has
\[H_X=H_Z=\begin{pmatrix}1&1&1&1\end{pmatrix},\]
\[G_X=\begin{pmatrix}1&1&0&0\\1&0&1&0\end{pmatrix},\quad G_Z=\begin{pmatrix}1&0&1&0\\1&1&0&0\end{pmatrix}.\]
All logical rows have even weight and hence commute with the checks. The first $X$ row and first $Z$ row overlap on one qubit and anticommute; different logical pairs overlap on an even number and commute. Direct multiplication gives $G_XG_Z^T=I_2$. We use the same quotient-and-pairing procedure for the BB codes.

The distinction from Section III is the matrix axis: $H_X,H_Z,G_X,G_Z$ have physical-qubit columns, whereas $D,L$ have circuit-fault columns. The code-level logical supports define the measured logical observables. Propagating each circuit fault and recording the observable flips then constructs the corresponding column of $L$.

\section{Completion and likelihood examples}\label{app:examples}
\subsection{Completion after a partial cut}
For prefix closure, take
\[D=\begin{pmatrix}1&1&0\\0&1&1\\0&1&0\end{pmatrix},\quad s=(1,1,1)^T=D(0,1,0)^T.\]
Its rows are the equations $x_1+x_2=1$, $x_2+x_3=1$, and $x_2=1$; its columns record which equations a fault toggles. With one row per block, columns one and two are owned by block one, and column three by block two. A first commitment $(x_1,x_2)=(1,0)$ closes row one; $x_3=1$ closes row two. Row three is then unsatisfied with no remaining variable. Earlier rows stay closed, but no completion exists. This is why the reference implementation declares a local failure and includes the same-engine global fallback in its measured cost.

\subsection{Configuration likelihood and logical-class weight}
For independent fault probabilities, logical-class weight is
\[Z_\lambda(s)=\sum_{x:Dx=s,\ Lx=\lambda}\exp[-\cost(x)].\]
The repair objective selects configurations, while the recorded logical outcomes evaluate the complete decoder. For logical degeneracy, take $D=(1\ 1\ 1)$, $L=(1\ 0\ 0)$, $s=1$, and independent probabilities $(.3,.2,.2)$. The detector row checks the parity $x_1+x_2+x_3$, while the logical row records $x_1$. Thus columns two and three have the same detector and logical signatures, and column one has the same detector signature but a different logical action. The syndrome-compatible configurations $100,010,001,111$ have masses $.192,.112,.112,.012$. Replacing $010$ by $100$ improves configuration likelihood, but changes from logical class zero, of total mass $.224$, to class one, of mass $.204$. A lower-cost representative can therefore be in the less probable logical class. This is an algebraic counterexample, not another simulated code experiment.

\section{Paired tests and uncertainty}\label{app:statistics}
For $f$ failures in $N$ independent samples, $\hat P=f/N$ and $z_W=1.9599639845$, the marginal Wilson endpoints are
\[\frac{\hat P+z_W^2/(2N)\ \pm\ z_W\sqrt{\hat P(1-\hat P)/N+z_W^2/(4N^2)}}{1+z_W^2/N}.\]
For paired discordances $n_{10},n_{01}$ and $M_{\rm disc}=n_{10}+n_{01}>0$,
\[p_{\rm exact}=\min\!\left\{1,\ 2\sum_{j=0}^{\min(n_{10},n_{01})}\binom{M_{\rm disc}}j2^{-M_{\rm disc}}\right\}.\]
Set $p_{\rm exact}=1$ for $M_{\rm disc}=0$. Holm adjustment sorts the 135 planned $p$-values and sets
\[\widetilde p_{(i)}=\min\{1,\max_{j\le i}(136-j)p_{(j)}\}.\]
Let $[a_{10},b_{10}]$ and $[a_{01},b_{01}]$ be 97.5\% Clopper--Pearson intervals for the two marginal discordance probabilities. The paired-effect interval $[a_{10}-b_{01},b_{10}-a_{01}]$ has at least 95\% coverage by the union bound, including zero-discordance cases. It is conservative and does not replace the exact paired test.

\section{Evaluation grid}\label{app:grid}
Table~\ref{tab:failgrid} reports the original observed failure counts. R1 and R2 denote repair at radii one and two. Radius-one BF repair lowers the observed failure count relative to direct cutting in 13 of the 15 configurations and ties in the other two. The reductions span both code families. For example, repair reduces failures from 55 to 39 on 1,024 paired shots for the four-block distance-five surface-code circuit at $p=0.003$, and from 44 to 32 on 128 paired shots for BB72 at $p=0.015$. The magnitude of the improvement varies across configurations, so the counts are best read alongside the shot totals.

Increasing the repair radius has mixed effects. Compared with radius-one BF repair, radius-two repair yields fewer failures in five configurations, ties in five, and yields more failures in five. The comparison with joint decoding also depends on the circuit. At $p=0.006$, radius-one BF repair gives 45 failures versus 50 for joint BF in the distance-five circuit, while the distance-seven circuit gives 22 versus 17. These counts describe the observed outcomes. The paired tests account for which individual shots each method fixes or turns into a failure.

\begin{table*}[t]\centering\scriptsize
\caption{Logical failure counts on the complete evaluation grid.}\label{tab:failgrid}
\begin{tabular}{lrrrrrrrrrrrr}\toprule
&&&&\multicolumn{4}{c}{BF}&\multicolumn{2}{c}{LSD2}&\multicolumn{2}{c}{OSD0}\\
Code&$p$&Blocks&$N$&Joint&Cut&R1&R2&Joint&R1&Joint&R1\\\midrule
Surface $d=3$&.001&4&512&10&7&5&8&8&5&10&5\\
Surface $d=3$&.003&4&512&30&29&25&25&28&24&30&25\\
Surface $d=3$&.006&4&512&59&63&54&56&51&51&59&53\\
Surface $d=5$&.001&4&256&1&2&2&1&0&1&1&2\\
Surface $d=5$&.003&4&1024&46&55&39&36&38&33&46&38\\
Surface $d=5$&.006&4&256&50&65&45&47&47&38&50&43\\
Surface $d=7$&.001&4&64&0&1&0&0&0&0&0&0\\
Surface $d=7$&.003&4&64&2&3&0&0&2&0&2&0\\
Surface $d=7$&.006&4&64&17&29&22&17&16&21&17&24\\
BB72&.005&4&128&6&9&5&5&6&4&6&5\\
BB72&.015&4&128&47&44&32&30&47&30&47&32\\
BB144&.005&4&128&0&0&0&0&0&0&0&0\\
BB144&.015&4&128&12&13&10&7&10&6&12&8\\
Surface $d=3$&.003&8&256&26&26&18&21&23&21&26&18\\
Surface $d=5$&.003&8&64&9&9&3&4&8&4&9&3\\\bottomrule
\end{tabular}\end{table*}
